\documentclass[11pt]{article}
\input{../../notes-he/structure}
\hypersetup{pdftitle={},pdfauthor={}}
\begin{document}
\begin{center}
{\LARGE\sffamily\bfseries\color{deepblue} קוויז בונוס --- שבוע 5}\\[.2cm]
{\large פונקציות על רשימות, ללא \eng{LinearAlgebra}}
\end{center}

\noindent שתי השאלות הן שאלות בונוס אופציונליות. בסביבת הבדיקה אין את החבילה
\eng{\texttt{LinearAlgebra}}, ולכן אין \eng{\texttt{dot}} או \eng{\texttt{norm}}:
יש לכתוב את הלולאה בעצמכם.

\section*{שאלה 2 --- מכפלה פנימית מנורמלת}
כתבו פונקציה המקבלת שתי רשימות של מספרים ומחזירה את המכפלה הפנימית שלהן,
מנורמלת לאורכים:
\[
\frac{a\cdot b}{|a|\,|b|}.
\]
חתימת הפונקציה:
\begin{latin}
\begin{lstlisting}
function norm_inner_product(v1, v2)
    # ...
end
\end{lstlisting}
\end{latin}
פלט הפונקציה: מספר. ניתן להניח שהקלט הוא אכן שתי רשימות המכילות מספרים מסוג
\eng{\texttt{Int}} או \eng{\texttt{Float64}} בלבד.
\begin{itemize}
\item אם הרשימות אינן באותו אורך, הפונקציה צריכה להחזיר \eng{\texttt{nothing}}.
\item אם הרשימות ריקות, הפונקציה צריכה להחזיר \eng{\texttt{0}}.
\end{itemize}
וקטור האפס הופך את המכפלה המנורמלת לבלתי מוגדרת; הבדיקות אינן כוללות מקרה כזה.

דוגמה: עבור \eng{\texttt{v1 = [3.0, 4.0]}} ו-\eng{\texttt{v2 = [6.0, 8.0]}} התוצאה
היא \eng{\texttt{1.0}}. עבור \eng{\texttt{[1.0, 0.0]}} ו-\eng{\texttt{[0.0, 1.0]}}
התוצאה היא \eng{\texttt{0.0}}.

\section*{שאלה 4 --- ספירת זוגות אורתוגונליים}
שני וקטורים אורתוגונליים זה לזה אם המכפלה הפנימית ביניהם היא אפס. כתבו פונקציה
המקבלת רשימה של וקטורים (רשימה של רשימות) ומחזירה את מספר הזוגות של הווקטורים
הניצבים זה לזה:
\begin{latin}
\begin{lstlisting}
function num_of_orthogonal(vectors)
    # ...
end
\end{lstlisting}
\end{latin}
פלט הפונקציה: מספר הזוגות של הווקטורים הניצבים זה לזה ברשימה.

\begin{latin}
\begin{lstlisting}
num_of_orthogonal([[1, 0, 0], [0, 1, 0], [0, 0, 1]])              # 3
num_of_orthogonal([[0, 0, 0], [0, 1, 0], [0, 0, 1]])              # 3
num_of_orthogonal([[0, 0], [1, 2], [10, 5]])                      # 2
num_of_orthogonal([[1, 1, 1, 1], [2, 1, 3, 3], [0, 0, 100, 33],
                   [8, 8, 8, 1.5], [9, 9, 9, 9]])                 # 0
\end{lstlisting}
\end{latin}
ניתן להניח כי הקלט הוא אכן רשימה, שהרשימות הפנימיות באותו אורך, ושאינן מכילות
ערכים שאינם מספרים. כל זוג נספר פעם אחת בלבד: הזוג \eng{(i, j)} והזוג
\eng{(j, i)} הם אותו זוג, ואין לספור וקטור עם עצמו. המנעו מכפילויות קוד;
היעזרו בפונקציית עזר אחת למכפלה פנימית בשתי השאלות.

\begin{center}
\begin{tabular}{lc}
\toprule
קריאה & תוצאה \\
\midrule
\eng{\texttt{norm\_inner\_product([1.0, 0.0], [-1.0, 0.0])}} & \eng{\texttt{-1.0}} \\
\eng{\texttt{norm\_inner\_product([1, 2, 3], [1, 2])}} & \eng{\texttt{nothing}} \\
\eng{\texttt{norm\_inner\_product(Float64[], Float64[])}} & \eng{\texttt{0}} \\
\eng{\texttt{num\_of\_orthogonal}} על בסיס אורתונורמלי ב-4 ממדים & \eng{\texttt{6}} \\
\bottomrule
\end{tabular}
\end{center}
\end{document}
